\section{计划栈}

出于教学目的，把计划拆成若干层会更清楚：
\begin{enumerate}[nosep]
  \item \textbf{结果层：} 想要达到的终局状态是什么？
  \item \textbf{会话层：} 哪一段具体会话能推动这个结果？
  \item \textbf{入口层：} 哪个确切的第一步动作会启动这段会话？
  \item \textbf{环境层：} 什么样的物理布置能够降低进入摩擦？
  \item \textbf{触发层：} 预期中的决策窗口会在何时、何地出现？
\end{enumerate}

\begin{proposition}[执行选择压缩]
如果一个计划在决策窗口到来之前解决了 $k \geq 1$ 个未来微决策，那么该窗口中的执行选择复杂度会从 $D$ 降到至多 $D-k$。
\end{proposition}

\begin{proof}
按照构造，每解决一个微决策，就少了一项必须在执行时刻再做出的行动相关决定。因此，若在窗口前解决了 $k$ 项这样的内容，剩余决策数至少减少 $k$，于是未来的决策复杂度至多为 $D-k$。
\end{proof}

\begin{corollary}[有限定的二元执行优势]
如果一个计划把执行复杂度压缩为一个二元选择---执行预先规定的动作，或者拒绝它---那么与任何其他在别的方面相似、但执行复杂度严格更大的计划相比，该计划在行动当下的执行负担都将弱更低。
\end{corollary}

\begin{proof}
一个二元选择的执行复杂度为 $D=1$。任何其他在别的方面相似、但复杂度严格更大的计划，都有 $D>1$。根据决策负载假设，更大的执行复杂度会使行动当下的负担弱增加。因此，二元选择计划会使这一负担弱降低。
\end{proof}

\begin{discussion}
这个结果比旧版表述更窄，也更诚实。它并不主张二元框架能够保证执行。它只是在说：在这个模型内部，减少那些仍需实时决定的问题数量，就会减少必须在窗口里支付的决策负担。
\end{discussion}